%% kapa: exam=1 round=37 year=2026 period=2 form=A subject=appraisal_law subject_ko=감정평가관계법규 seq=22 number=22
%% source: Q-Net 감정평가사 기출문제 (cst00309) artlSeq=5258924 file=raw/1008/5258924/2026년 제37회 감정평가사 1차 2교시.pdf
%% conversion: kordoc 4.24.51 -> tools/gen1.py (auto); review=proofread-sonnet; flags=none
\begin{question}{22}
국유재산법령상 변상금의 징수에 관한 조문의 일부이다. ( )에 들어갈 숫자로 옳은 것은?
\begin{framedtext}
○ 중앙관서의 장등은 무단점유자에 대하여 대통령령으로 정하는 바에 따라 그 재산에 대한 사용료나 대부료의 100분의 ( ㄱ )에 상당하는 변상금을 징수한다. \par
○ 변상금은 무단점유를 하게 된 경위, 무단점유지의 용도 및 해당 무단점유자의 경제적 사정 등을 고려하여 대통령령으로 정하는 바에 따라 ( ㄴ )년의 범위에서 징수를 미루거나 나누어 내게 할 수 있다. \par
\end{framedtext}
\begin{choices}
  \item ㄱ: 110, ㄴ: 3
  \item ㄱ: 110, ㄴ: 5
  \item ㄱ: 120, ㄴ: 3
  \item ㄱ: 120, ㄴ: 5
  \item ㄱ: 150, ㄴ: 5
\end{choices}
\end{question}
\answer{4}
